% ═══════════════════════════════════════════════════════════════
% LERNBLATT: Repaso - La ropa y los adjetivos
% Spanisch Klasse 8 - Max. 2 A4-Seiten
% ═══════════════════════════════════════════════════════════════
% DESIGN-PRINZIPIEN:
% - Sans-Serif (Helvetica): Fließtext
% - Serif (Palatino): Spanische Vokabeln (visuelle Distinktion)
% - Minimale Rahmen, kein "visual noise"
% ═══════════════════════════════════════════════════════════════

\documentclass[11pt, a4paper]{article}

\usepackage[utf8]{inputenc}
\usepackage[T1]{fontenc}
\usepackage[ngerman]{babel}

% --- SCHRIFTEN ---
\usepackage{helvet}
\renewcommand{\familydefault}{\sfdefault}
\usepackage{palatino}
\newcommand{\seriftext}[1]{{\fontfamily{ppl}\selectfont #1}}

\usepackage{microtype}
\usepackage[margin=1.8cm, top=2cm, bottom=1.5cm]{geometry}
\usepackage{enumitem}
\usepackage{tabularx}
\usepackage{booktabs}
\usepackage{array}
\usepackage{multicol}

\usepackage{fancyhdr}
\usepackage{lastpage}

% ═══════════════════════════════════════════════════════════════
% FARBEN
% ═══════════════════════════════════════════════════════════════

\usepackage{xcolor}

% Strukturfarben
\definecolor{hellgrau}{HTML}{F5F5F5}
\definecolor{rahmen}{HTML}{CCCCCC}
\definecolor{akzent}{HTML}{666666}
\definecolor{spanisch}{HTML}{c41e3a}

% ═══════════════════════════════════════════════════════════════
% BOXEN
% ═══════════════════════════════════════════════════════════════

\usepackage{tcolorbox}
\tcbuselibrary{skins}

\newtcolorbox{lernzielbox}{
    colback=hellgrau, colframe=hellgrau, boxrule=0pt, arc=0pt,
    left=10pt, right=10pt, top=8pt, bottom=8pt,
    before skip=6pt, after skip=8pt
}

\newtcolorbox{grammatikbox}[1][]{
    colback=white, colframe=white, boxrule=0pt, arc=0pt,
    left=10pt, right=6pt, top=6pt, bottom=6pt,
    borderline west={3pt}{0pt}{akzent},
    before skip=4pt, after skip=4pt
}

\newtcolorbox{merkbox}[1][]{
    colback=hellgrau, colframe=hellgrau, boxrule=0pt, arc=0pt,
    left=10pt, right=10pt, top=8pt, bottom=8pt,
    borderline north={2pt}{0pt}{spanisch},
    before skip=6pt, after skip=6pt
}

% ═══════════════════════════════════════════════════════════════
% BEFEHLE
% ═══════════════════════════════════════════════════════════════

\newcommand{\es}[1]{\seriftext{\textbf{#1}}}
\newcommand{\de}[1]{\textit{\textcolor{akzent}{#1}}}
\newcommand{\gram}[1]{\textsc{#1}}
\newcommand{\abmark}[1]{{\footnotesize\textcolor{akzent}{$\rightarrow$ AB #1}}}

\setlength{\parindent}{0pt}
\setlength{\parskip}{5pt}

% ═══════════════════════════════════════════════════════════════
% KOPF- UND FUSSZEILE
% ═══════════════════════════════════════════════════════════════

\pagestyle{fancy}
\fancyhf{}
\renewcommand{\headrulewidth}{0.5pt}
\renewcommand{\footrulewidth}{0pt}
\lhead{\footnotesize\textsc{Spanisch} · Klasse 8}
\rhead{\footnotesize La ropa y los adjetivos}
\cfoot{\footnotesize\thepage\,/\,\pageref{LastPage}}

% ═══════════════════════════════════════════════════════════════
% DOKUMENT
% ═══════════════════════════════════════════════════════════════

\begin{document}

% --- TITEL ---
\begin{center}
    {\LARGE\bfseries La ropa y los adjetivos}\\[3pt]
    {\small Kleidung beschreiben · Adjektive · Genus-Anpassung}
\end{center}

\vspace{4pt}

% --- EINLEITUNG ---
Wie beschreibst du dein Lieblingsoutfit auf Spanisch?
Mit den richtigen Adjektiven wird aus \es{un jersey} ein \es{jersey bonito y cómodo}!

% --- LERNZIELE ---
\begin{lernzielbox}
\textbf{Das lernst du:}\quad
Kleidungsstücke benennen · Adjektive zur Beschreibung · Genus-Anpassung (-o/-a/-e) · Redemittel ``Über Kleidung sprechen''
\end{lernzielbox}

% ═══════════════════════════════════════════════════════════════
% ZWEI SPALTEN
% ═══════════════════════════════════════════════════════════════

\begin{multicols}{2}

% ---------------------------------------------------------------
% 1. KLEIDUNG
% ---------------------------------------------------------------
\textbf{\large 1. La ropa} \de{-- Kleidung}

\vspace{4pt}

\begin{tabular}{@{}ll@{}}
\es{la camiseta} & \de{das T-Shirt} \\[0.4em]
\es{la camisa} & \de{das Hemd} \\[0.4em]
\es{el jersey} & \de{der Pullover} \\[0.4em]
\es{la sudadera} & \de{der Kapuzenpulli} \\[0.4em]
\es{los pantalones} & \de{die Hose} \\[0.4em]
\es{el vestido} & \de{das Kleid} \\[0.4em]
\es{la falda} & \de{der Rock} \\[0.4em]
\es{la chaqueta} & \de{die Jacke} \\[0.4em]
\es{el abrigo} & \de{der Mantel} \\[0.4em]
\es{los zapatos} & \de{die Schuhe} \\[0.4em]
\es{las botas} & \de{die Stiefel} \\[0.4em]
\es{las zapatillas} & \de{die Turnschuhe} \\
\end{tabular}

\vspace{4pt}
{\footnotesize\textit{Tipp:} \es{Los pantalones} und \es{los zapatos} stehen im Spanischen immer im Plural!}

\abmark{Aufgabe 1}

\vspace{8pt}

% ---------------------------------------------------------------
% 2. ADJEKTIVE
% ---------------------------------------------------------------
\textbf{\large 2. Los adjetivos} \de{-- Adjektive}

\vspace{4pt}

\begin{tabular}{@{}ll@{}}
\es{bonito/a} & \de{hübsch} \\[0.4em]
\es{moderno/a} & \de{modern} \\[0.4em]
\es{cómodo/a} & \de{bequem} \\[0.4em]
\es{ancho/a} & \de{weit} \\[0.4em]
\es{estrecho/a} & \de{eng} \\[0.4em]
\es{largo/a} & \de{lang} \\[0.4em]
\es{corto/a} & \de{kurz} \\[0.4em]
\es{deportivo/a} & \de{sportlich} \\[0.4em]
\es{sencillo/a} & \de{schlicht} \\[0.4em]
\es{elegante} & \de{elegant} \\
\end{tabular}

\vspace{4pt}
{\footnotesize\textit{Merke:} \es{elegante} hat nur eine Form -- keine Anpassung nötig!}

\abmark{Aufgabe 2}

\columnbreak

% ---------------------------------------------------------------
% 3. GRAMMATIK: GENUS-ANPASSUNG
% ---------------------------------------------------------------
\textbf{\large 3. Adjektive anpassen}

\vspace{4pt}

Im Spanischen steht das Adjektiv \textbf{nach} dem Nomen und passt sich im \textbf{Genus} an:

\begin{grammatikbox}
\textbf{Regel 1: Adjektive auf -o/-a}

Das Adjektiv ``kopiert'' das Genus des Nomens:

\vspace{3pt}
{\small
\begin{tabular}{@{}ll@{}}
\es{el jersey bonit\textbf{o}} & (maskulin) \\[0.3em]
\es{la camiseta bonit\textbf{a}} & (feminin) \\[0.3em]
\es{los zapatos cómod\textbf{os}} & (mask. Plural) \\[0.3em]
\es{las botas cómod\textbf{as}} & (fem. Plural) \\
\end{tabular}}
\end{grammatikbox}

\vspace{4pt}

\begin{grammatikbox}
\textbf{Regel 2: Adjektive auf -e oder Konsonant}

Diese Adjektive \textbf{ändern sich nicht}:

\vspace{3pt}
{\small
\begin{tabular}{@{}ll@{}}
\es{el vestido elegant\textbf{e}} & \\[0.3em]
\es{la falda elegant\textbf{e}} & (gleich!) \\
\end{tabular}}
\end{grammatikbox}

\vspace{4pt}
{\footnotesize\textit{Eselsbrücke:} ``Das Adjektiv ist der Schatten des Nomens -- es folgt ihm überall!''}

\abmark{Aufgabe 2 + 3}

\vspace{8pt}

% ---------------------------------------------------------------
% 4. REDEMITTEL
% ---------------------------------------------------------------
\textbf{\large 4. Über Kleidung sprechen}

\vspace{4pt}

\textbf{Beschreiben:}\\
{\small
\es{El/La ... es muy ...}\\
\de{= Der/Die/Das ... ist sehr ...}

\vspace{3pt}
\es{Llevo un/una ... ...}\\
\de{= Ich trage einen/eine ... ...}}

\vspace{6pt}

\textbf{Bewerten:}\\
{\small
\es{Me gusta el/la ... porque es ...}\\
\de{= Mir gefällt der/die ... weil er/sie ... ist}

\vspace{3pt}
\es{El/La ... es demasiado caro/a.}\\
\de{= Der/Die ... ist zu teuer.}}

\vspace{6pt}

\textbf{Beispiele:}\\
{\footnotesize
\es{La camiseta es muy cómoda.}\\
\es{Me gusta el jersey porque es moderno.}\\
\es{Llevo una falda elegante.}}

\abmark{Aufgabe 3 + 4}

\end{multicols}

% ═══════════════════════════════════════════════════════════════
% MERKKASTEN
% ═══════════════════════════════════════════════════════════════

\begin{merkbox}
\textbf{Zusammenfassung: Adjektiv-Anpassung}

\vspace{4pt}
\begin{center}
\begin{tabular}{|l|c|c|c|c|}
\hline
\textbf{Adjektiv} & \textbf{mask. Sg.} & \textbf{fem. Sg.} & \textbf{mask. Pl.} & \textbf{fem. Pl.} \\
\hline
bonito & bonit\textbf{o} & bonit\textbf{a} & bonit\textbf{os} & bonit\textbf{as} \\
\hline
cómodo & cómod\textbf{o} & cómod\textbf{a} & cómod\textbf{os} & cómod\textbf{as} \\
\hline
elegante & elegante & elegante & elegantes & elegantes \\
\hline
\end{tabular}
\end{center}

\vspace{4pt}
{\small\textbf{Kurz:} -o/-a passt sich an · -e bleibt gleich · Plural: +s (nach Vokal) oder +es (nach Konsonant)}
\end{merkbox}

\vspace{6pt}

% ═══════════════════════════════════════════════════════════════
% PARTNERÜBUNG
% ═══════════════════════════════════════════════════════════════

\begin{tcolorbox}[
    colback=white, colframe=rahmen, boxrule=0.5pt, arc=0pt,
    left=8pt, right=8pt, top=6pt, bottom=6pt,
    borderline north={1.5pt}{0pt}{spanisch}
]
{\small\textbf{Partnerübung: Beschreibe dein Outfit!}}

\vspace{4pt}
\footnotesize
Beschreibe deinem Partner 2--3 Kleidungsstücke, die du heute trägst. Verwende die Redemittel!

\vspace{4pt}
\textbf{Beispiel:}\\
A: \es{Hoy llevo una camiseta blanca. Es muy cómoda.}\\
B: \es{Yo llevo un jersey azul. Es moderno y bonito.}

\vspace{4pt}
\textit{Tipp: Achte auf die Adjektiv-Endung!}
\end{tcolorbox}

\end{document}
